\chapter{Гипотеза GC: Green-контактная оценка}
\noindent В главе рассматривается открытая гипотеза GC о Carleson-контроле полной массы контактных ячеек. Для множества Мандельброта строятся суммируемые блоковые масштабы и оценивается близость интерфейсов к предельному множеству; двойной подсчёт и оценки низкой Green-высоты дают строгие редукции к контактной массе. Недоказанным остаётся её равномерный контроль с требуемым множителем радиуса.

\section{Открытый Carleson-бюджет}

\begin{openestimate}[Carleson-бюджет для Green-интерфейсов (\(\mathsf{GC}\))]
\label{thm:green-contact-carleson}
Пусть \(X\) — компактный полный континуум в плоскости, \(G\) —
непрерывная функция Грина, продолженная нулём на \(X\), а
\(E_s:=\{z:G(z)\le s\}\). Рассмотрим вложенную компактную фильтрацию
\(Y_n\searrow X\), заданную ограничениями
\[
 Y_n=Y_{n-1}\cap P_n^{-1}(\overline{\mathbb D}),
\]
где каждый \(P_n\) не константен и голоморфен в окрестности \(Y_0\).
Предположим, что
\[
 \Sigma_n\subseteq
 Y_{n-1}\cap P_n^{-1}(\partial\mathbb D)\cap E_{g_n},
\]
 где каждое \(\Sigma_n\) компактно и его интерфейсный граф конечен; все
 контакты этого графа предполагаются включёнными в \(\Sigma_n\), как
 описано выше. Пусть
контактные ячейки блока лежат в \(\mu_m\)-окрестности \(X\), где
\(g_n,\mu_m\ge0\). Для
фиксированных \(x\in X\), \(r>0\) используем радиусы \(R_m\),
определённые выше, предполагаем \(R_m>0\) при всех достаточно больших
\(m\), и положим
\[
 \epsilon_m:=
 \sqrt{\frac{\mu_m}{r}}+
 \sum_{\lambda_m<n\le\lambda_{m+1}}\sqrt{g_n}.
\]
При \(\sum_m\epsilon_m<\infty\) требуется доказать существование
\(B_{x,r}>0\) и \(m_0\), не зависящих от \(m\), таких, что для всех
\(m\ge m_0\)
\begin{equation}
 \Xi_m(\C)\le B_{x,r} R_m\epsilon_m.
 \label{eq:green-contact-carleson}
\end{equation}
\end{openestimate}
Это открытое утверждение не следует из одной лишь близости интерфейсов
к \(X\): оценка \(\mu_m\to0\) не контролирует кратность встреч ячеек с
разными \(\Sigma_n\). Для доказательства нужен локальный
Carleson-контроль контактной меры, включая кратности критических карт
и перекрытия обратных ветвей. Утверждение является открытой целью, а
не доказанным следствием одних голоморфных ограничений. Если широкий
класс фильтраций в формулировке окажется слишком общим, требуется
выделить независимо проверяемые условия на локальные степени и
перекрытия, сохранив применимость к различным голоморфным
Green-фильтрациям, а не только к одной рекурсии.

\section{Суммируемые масштабы для множества Мандельброта}

\begin{proposition}[Суммируемые блоковые масштабы для \(\M\)]
\label{prop:mandelbrot-green-block-scales}
Для \(c\in\partial\M\), \(r>0\) и \(n\ge1\) положим
\begin{equation}
 \Sigma_n(c,r):=C_{n-1}(c,r)\cap\{z:|p_n(z)|=2\}.
 \label{eq:green-separator}
\end{equation}
Существуют числа \(s_m>0\) и строго возрастающие целые
\(\lambda_m\ge m\), \(m\ge1\), такие, что при \(m\ge1\), если положить
\[
 g_n:=2^{2-n}\log3,\qquad
 \mu_m:=8\cdot2^{-m}+\kappa(s_m),\qquad
 \epsilon_m(r):=
 \sqrt{\frac{\mu_m}{r}}+
 \sum_{\lambda_m<n\le\lambda_{m+1}}\sqrt{g_n},
\]
выполняются \(\mu_m\le9\cdot2^{-m}\) при \(m\ge2\) и
\(\sum_{m\ge1}\epsilon_m(r)<\infty\) для каждого \(r>0\).
Последовательности \(s_m,\lambda_m\) можно выбрать один раз независимо
от \(c,r\).
Кроме того, если ячейка \(Q_{m,\alpha}\) сетки
\eqref{eq:dyadic-grid-cells} встречает интерфейс \(\Sigma_n\) при
\(m\ge2\) и \(\lambda_m<n\le\lambda_{m+1}\), то
\[
 \sup_{z\in Q_{m,\alpha}}\dist(z,\M)\le\mu_m.
\]
\end{proposition}
\begin{proof}
По теореме~\ref{thm:green-projection} имеем
\(\kappa(s)=\dH(E_s,\M)\to0\) при \(s\downarrow0\). Положим
\(s_1=1\), \(\lambda_1=1\). Для \(m\ge2\) выбираем \(s_m>0\) так,
чтобы
\[
 \kappa(s_m)\le2^{-m},\qquad s_m\le2^{-m},\qquad
 s_m\le s_{m-1}/2.
\]
Затем рекурсивно выбираем \(\lambda_m>\lambda_{m-1}\), \(\lambda_m\ge m\),
настолько большим, что \(2^{1-\lambda_m}\log3\le s_m\). Для
\(m\ge2\), \(\lambda_m<n\le\lambda_{m+1}\) и
\(\zeta\in\Sigma_n(c,r)\) имеем \(\zeta\in O_{n-1}\); по
лемме~\ref{lem:green-depth}
\[
 G(\zeta)\le2^{2-n}\log3=g_n
 \le2^{1-\lambda_m}\log3\le s_m.
\]
Значит, \(\dist(\zeta,\M)\le\kappa(s_m)\). Если
\(Q_{m,\alpha}\cap\Sigma_n(c,r)\ne\varnothing\), выбираем
\(\zeta\in Q_{m,\alpha}\cap\Sigma_n(c,r)\). Оценка диаметра ячейки
\(\operatorname{diam}Q_{m,\alpha}\le8\cdot2^{-m}\) даёт заявленный
контроль расстояния до \(\M\).

Наконец,
\[
 \sum_{n>\lambda_m}\sqrt{g_n}
 \le\frac{\sqrt{2\log3}}{1-2^{-1/2}}\,2^{-\lambda_m/2}
 \le\frac{\sqrt{2\log3}}{1-2^{-1/2}}\,2^{-m/2},
\]
а \(\sqrt{\mu_m/r}\le3r^{-1/2}2^{-m/2}\). Оба вклада суммируемы
по \(m\ge2\); начальный член \(\epsilon_1(r)\) конечен. Это доказывает
утверждение.
\end{proof}

\section{Редукции к гипотезе GC: упаковка контактов}
\label{sec:gc-open-reductions}

Зафиксируем $c\in\partial\M$, $r>0$ и будем писать
$\Sigma_n=\Sigma_n(c,r)$.
Для сетки $Q_{m,\alpha}$ положим
\[
 I_{n,m}:=\{\alpha:Q_{m,\alpha}\cap\Sigma_n(c,r)\ne\varnothing\},\qquad
 v_{m,\alpha}:=\#\{n:\lambda_m<n\le\lambda_{m+1},\
                         \alpha\in I_{n,m}\}.
\]
При выборе центров $z_{m,\alpha}\in Q_{m,\alpha}$ контактная мера
имеет вид
\[
 \Xi_m:=\sum_{\alpha\in J_m^2}
 v_{m,\alpha}\operatorname{diam}(Q_{m,\alpha})\delta_{z_{m,\alpha}}.
\]
Поскольку все ячейки, встречающие $\Sigma_n$, покрывают интерфейс,
полный набор этих ячеек является допустимым контактным разрезом на
каждом уровне. Поэтому
\begin{equation}
 \mathfrak c_m\le
 \sum_{\lambda_m<n\le\lambda_{m+1}}
       \sum_{\alpha\in I_{n,m}}\operatorname{diam}Q_{m,\alpha}
 =\sum_{\alpha\in J_m^2}
       v_{m,\alpha}\operatorname{diam}Q_{m,\alpha}
 =\Xi_m(\C).
 \label{eq:gc-mass-reduction}
\end{equation}
Равенство в середине — точный двойной подсчёт; оно учитывает кратность
встреч одной ячейки с разными интерфейсами. Оценка диаметра даёт лишь
\[
 \mathfrak c_m\le8\cdot2^{-m}\sum_{\alpha\in J_m^2}v_{m,\alpha},
 \qquad v_{m,\alpha}\le\lambda_{m+1}-\lambda_m,
 \qquad |J_m^2|=4^m.
\]
Эти конечные оценки не дают множителя $R_m$ и потому не влекут
Carleson-бюджет $\mathsf{GC}$.

Предложение~\ref{prop:mandelbrot-green-block-scales} даёт точную
локализацию контактов: если $Q_{m,\alpha}$ встречает $\Sigma_n$ в
блоке $\lambda_m<n\le\lambda_{m+1}$, то
\[
 \sup_{z\in Q_{m,\alpha}}\dist(z,\M)\le\mu_m\le9\cdot2^{-m}.
\]
Кроме того, для $\zeta\in\Sigma_n$ выполняется точная нормировка
\[
 g_\zeta(p_n(\zeta))=2^{n-1}G(\zeta).
\]
Она позволяет разбить контакты по динамической высоте. Выбрав
$0<\tau_m\le2^{-m}$ так, чтобы
$\kappa(2^{-\lambda_m}\tau_m)\le2^{-2m}$, и положив
\[
 \Sigma_n^-(\tau_m):=
 \{\zeta\in\Sigma_n:g_\zeta(p_n(\zeta))\le\tau_m\},
 \qquad
 a_{n,m}^-:=\sup\bigl(\{G(\zeta):\zeta\in\Sigma_n^-(\tau_m)\}\cup\{0\}\bigr),
\]
получаем для $\lambda_m<n\le\lambda_{m+1}$
\begin{equation}
 \forall\zeta\in\Sigma_n^-(\tau_m),\quad
 \dist(\zeta,\M)\le2^{-2m},\qquad
 \sum_{\lambda_m<n\le\lambda_{m+1}}\sqrt{a_{n,m}^-}
 \le\frac{2^{-m}}{1-2^{-1/2}}.
 \label{eq:gc-low-height-budget}
\end{equation}
Выбор $\tau_m$ возможен ввиду $\kappa(s)\to0$ при $s\downarrow0$.
Действительно, для таких $\zeta$
\[
 G(\zeta)=2^{1-n}g_\zeta(p_n(\zeta))
 \le2^{1-n}\tau_m\le2^{-\lambda_m}\tau_m,
\]
что даёт оценку расстояния по выбору $\tau_m$. Кроме того,
\[
 \sum_{\lambda_m<n\le\lambda_{m+1}}\sqrt{a_{n,m}^-}
 \le\sqrt{\tau_m}\sum_{n>\lambda_m}2^{(1-n)/2}
 =\frac{\sqrt{\tau_m}\,2^{-\lambda_m/2}}{1-2^{-1/2}}
 \le\frac{2^{-m}}{1-2^{-1/2}}.
\]
Эта строгая оценка контролирует лишь низковысотную часть: она не
ограничивает число контактных ячеек и не даёт сумму их кратностей.

Локальные обратные карты дают ещё один возможный источник весов, но
требуют равномерных оценок. Для фиксированных $n$ и $\zeta$ обозначим
через $q_n(\zeta)$ локальную степень $p_n$ в $\zeta$. Тогда
\[
 p_n(\zeta+w)-p_n(\zeta)=w^{q_n(\zeta)}h_{n,\zeta}(w),
 \qquad h_{n,\zeta}(0)\ne0.
\]
На достаточно малом диске можно выбрать $0<a\le|h_{n,\zeta}(w)|\le b$;
для $0<t\le a\delta^{q_n(\zeta)}$ отсюда следует
\begin{equation}
 \overline B\!\left(\zeta,(t/b)^{1/q_n(\zeta)}\right)
 \subseteq\{|z-\zeta|\le\delta,\ |p_n(z)-p_n(\zeta)|\le t\}
 \subseteq\overline B\!\left(\zeta,(t/a)^{1/q_n(\zeta)}\right).
 \label{eq:local-degree-pullback}
\end{equation}
При регулярном центре $q_n=1$; в критическом случае степень может быть
больше единицы. Для каждого фиксированного $n$ компактность $\Sigma_n$
даёт конечное покрытие такими картами, однако их число, радиусы,
константы и кратности перекрытий зависят от $n$. Глобальное тождество
следует из $\deg p_1=1$ и
$\deg p_{n+1}=2\deg p_n$, поскольку $p_{n+1}=p_n^2+c$:
\[
 \sum_{p_n'(\zeta)=0}(q_n(\zeta)-1)
 =\deg p_n'=2^{n-1}-1
\]
также не даёт локальной оценки на интерфейсе блока. Наконец, малость
$G(\zeta)$ и близость $\zeta$ к $\M$ сами по себе не дают малости
евклидова радиуса $t$ в образе $p_n$; для суммирования оценок
\eqref{eq:local-degree-pullback} нужны отдельные равномерные оценки
образных радиусов, степеней и перекрытий.

Итак, доказанные сведения дают точный переход от контактных ячеек к
полной массе $\Xi_m(\C)$ и суммируемую толщину их коллара. Открытым
остаётся именно локальный Carleson-контроль этой массы через
$R_m\epsilon_m$; ни близость интерфейсов к $\M$, ни конечность графов,
ни общий подсчёт критических кратностей его не доказывают.
